% tbsrc/resolution.tex — Detector Resolution and the Point-Spread Function
% Status: written (first draft). Blur sources; pixelated vs monolithic; the
% Gaussian PSF operator G; resolution injected at simulation time (not in the
% reconstruction), reconciling the forward model of Sec. projection/intro.
\section{Detector Resolution and the Point-Spread Function}
\label{sec:resolution}

The projector of Section~\ref{sec:projection} integrates infinitely thin lines:
it has no resolution of its own. Real PET data are blurred, and the previous
sections were careful to keep that blur out of the geometry. This section is
where it goes back in --- as a separate operator $\Gop$ --- and explains where,
in this tutorial, it enters the calculation.

\subsection{Sources of blur}
\label{sec:resolution:sources}

The spatial resolution of a PET scanner is limited by several physical effects,
none of which the reconstruction can undo by itself~\cite{moses2011}:
\begin{itemize}
  \item \emph{Positron range}: the positron travels a short distance before
        annihilating, so the annihilation is displaced from the emission ---
        sub-millimetre for $^{18}$F, larger for higher-energy emitters.
  \item \emph{Photon non-collinearity}: the two photons are not emitted exactly
        back to back (the pair carries residual momentum), spreading the angle
        by about $\pm0.25^\circ$; the resulting blur grows with the ring
        diameter, reaching a couple of millimetres in a human-scale scanner.
  \item \emph{Intrinsic positioning}: how well the detector localizes each
        photon interaction --- set by the crystal size for a pixelated detector,
        or by the readout and the position estimator for a monolithic one.
  \item \emph{Depth of interaction (DOI)}: photons enter the detector
        obliquely and interact at varying depth; without DOI information this
        smears the LOR, worsening off-centre (the parallax effect).
\end{itemize}
These contributions add approximately in quadrature to a total resolution,
quoted as a full width at half maximum ($\FWHM$). It is a property of the
scanner and the tracer, supplied to the model as data --- \code{RecoCrysp}
derives none of it.

\subsection{Pixelated versus monolithic detectors}
\label{sec:resolution:detectors}

The intrinsic term is where the two detector models of
Section~\ref{sec:geometry} differ sharply. A pixelated crystal of width $\ell$
localizes a photon only to within that crystal; a uniform uncertainty over a
width $\ell$ has standard deviation $\sigma=\ell/\sqrt{12}$, and the two
endpoints together give an intrinsic resolution of about $\ell/2$ FWHM at the
centre of the field of view. Resolution is thus tied to crystal size: finer
images demand smaller crystals, hence more readout channels.

A monolithic detector breaks this link. A large undivided block is read by a
photosensor array, and a neural network estimates the interaction position from
the light pattern; the accuracy is set by the light statistics and the network,
\emph{not} by the block size. CRYSP~\cite{crysp2026} reaches about
$3.5\,\mathrm{mm}$ FWHM with $50\,\mathrm{mm}$ blocks --- far finer than the
$\ell/2\approx25\,\mathrm{mm}$ a pixelated reading of the same block would give.
This decoupling is exactly why Section~\ref{sec:geometry:continuous} models the
detector as a continuous surface.

\subsection{The point-spread function as an operator}
\label{sec:resolution:psf}

Whatever its sources, the combined effect is that a point of activity is
recorded as a small blob: the system point-spread function (PSF). Over the
central field of view it is well approximated by an isotropic, spatially
invariant Gaussian of the total FWHM, which we write as a linear operator
$\Gop$ acting on the image --- convolution with that Gaussian,
\begin{equation}
  (\Gop\,\xv)(\bm u) = \int g(\bm u-\bm u')\,\xv(\bm u')\,\mathrm{d}\bm u' ,
  \qquad
  \sigma = \frac{\FWHM}{2\sqrt{2\ln 2}} \approx \frac{\FWHM}{2.355} ,
  \label{eq:res-gaussian}
\end{equation}
with $g$ the Gaussian of width $\sigma$. Like $\Aop$, $\Gop$ is applied without
forming a matrix; being a symmetric kernel it is self-adjoint, $\Gop=\Gop^{\mathsf T}$.
The resolution-aware forward operator is the composition $\Aop\,\Gop$, as
previewed in the full data model \eqref{eq:intro-full}. (The grid must sample
this Gaussian: this is the origin of the $\Delta\approx\tfrac13$--$\tfrac12\,\FWHM$
rule of Section~\ref{sec:grid:resolution}.) Treating the PSF as isotropic and
shift-invariant is an approximation --- real PSFs broaden towards the edge of
the field of view --- but an adequate one for the central-FOV examples here.

\subsection{Where the resolution enters}
\label{sec:resolution:where}

In this tutorial the resolution is injected \emph{at simulation time}, not in
the reconstruction. To generate realistic data for a known phantom
$\xv_{\mathrm{true}}$ we blur it and project the blurred image,
\begin{equation}
  \pred = \Aop\,(\Gop\,\xv_{\mathrm{true}}),
  \qquad
  \yv \sim \mathrm{Poisson}(\pred),
  \label{eq:res-sim}
\end{equation}
so the data $\yv$ carry the $3.5\,\mathrm{mm}$ blur. The reconstruction
(Section~\ref{sec:mlem}) then runs with the \emph{plain} projector $\Aop$.
Because the data are explained by $\Aop\,(\Gop\,\xv_{\mathrm{true}})$, the
iteration converges towards the blurred image $\Gop\,\xv_{\mathrm{true}}$: the
result is resolution-limited, resolving structures larger than the FWHM and not
those below it. This is the point of the resolution model --- one sees the limit
directly in the reconstructed image.

\begin{remark}
Equation~\eqref{eq:res-sim} is the \emph{image-space} way to inject resolution,
convenient because the phantom is known. A Monte-Carlo simulation does it the
\emph{detector} way instead --- jittering each LOR endpoint by the positioning
error, with no $\Gop$ operator at all. Both put the blur in the data, and the
reconstruction is identical. Building $\Gop$ \emph{into} the reconstruction
model (so that the iteration deconvolves towards $\xv_{\mathrm{true}}$) is
\emph{resolution recovery}: feasible with the same $\Gop$, but deliberately left
out of scope --- the examples reconstruct with $\Aop$ and report the
resolution-limited image.
\end{remark}
